% 代数部分
\usetikzlibrary{calc}

% 像空间
\newcommand\Img[0]{
    \mathrm{Img\,}
}
% 核：零空间
\newcommand\Ker[0]{
    \mathrm{Ker\,}
}
% 容：像空间的维数
\newcommand\rang[0]{
    \mathrm{rang\,}
}
% 佚：零空间的维数
\newcommand\nul[0]{
    \mathrm{null\,}
}
% 迹
\newcommand\ji[0]{\operatorname{tr}}
% 函数的限制：限制在定义域的某个子集中的函数
\newcommand\restr[2]{{% we make the whole thing an ordinary symbol
  \left.\kern-\nulldelimiterspace % automatically resize the bar with \right
  #1 % the function
  \vphantom{\big|} % pretend it's a little taller at normal size
  \right|_{#2} % this is the delimiter
  }
}
% 随附矩阵、附阵、伴随矩阵
\newcommand\co[0]{\operatorname{Co}}

% 交换集：交换元构成的集合
\newcommand\Cmu[1]{
    \mathrm{Comm\,\left(#1\right)}
}

% 商结构：两个代数结构的商结构

% ================= 可调参数（在导言区修改） =================
% 斜线与水平线的夹角（单位：度），默认 30 度
\newcommand{\quotientangle}{30}
% 斜线的总长度，默认 1.2em（可改为其他长度，如 1.5em、2ex 等）
\newcommand{\quotientlength}{2em}
% G 相对于斜线起点的水平偏移（负值向左移动），默认 -0.15em
\newcommand{\quotientGxshift}{-0.5em}
% G 相对于斜线起点的垂直偏移（正值向上移动），默认 0.2em
\newcommand{\quotientGyshift}{0.5em}
% H 相对于斜线终点的水平偏移（正值向右移动），默认 0.15em
\newcommand{\quotientHxshift}{0.5em}
% H 相对于斜线终点的垂直偏移（负值向下移动），默认 -0.2em
\newcommand{\quotientHyshift}{-0.5em}
% 斜线的线宽，默认使用当前 TikZ 线宽（\pgflinewidth）
% 如需加粗，可改为例如 0.8pt 或 1pt
\newcommand{\quotientlinewidth}{\pgflinewidth}
\newcommand{\quotientvshift}{0pt}
% =============================================================

% 商群命令：\quotient{G}{H}
% 在数学模式中使用，例如 $\quotient{G}{H}$
\newcommand{\quotient}[2]{%
  \mathord{% 保证作为普通数学符号的间距
    \tikz[baseline={([yshift=-\the\fontdimen22\textfont2 + \quotientvshift] mid)}]{
      % 定义斜线中点，位于坐标原点，作为对齐基准
      \coordinate (mid) at (0,0);
      % 计算斜线起点和终点（极坐标，关于中点对称）
      \coordinate (start) at ({180+\quotientangle}:{0.5*\quotientlength});
      \coordinate (end) at ({\quotientangle}:{0.5*\quotientlength});
      % 绘制斜线
      \draw[line width=\quotientlinewidth] (start) -- (end);
      % 放置 G，位置相对于中点偏移（默认在左上方）
      \node[inner sep=0pt] at ($(mid)+(\quotientGxshift,\quotientGyshift)$) {$#1$};
      % 放置 H，位置相对于中点偏移（默认在右下方）
      \node[inner sep=0pt] at ($(mid)+(\quotientHxshift,\quotientHyshift)$) {$#2$};
    }%
  }%
}

% 子群
\newcommand{\subg}{    
    \mathord{
        \tikz[baseline=0pt, trim left=-0.2em, trim right=0.8em] \draw (0,.35em) -- (.45em,0.6em) -- (.45em,0.1em)  -- cycle;
    }}
\newcommand{\subgeq}{    
    \mathord{
        \tikz[baseline=0pt, trim left=-0.2em, trim right=0.8em] \draw (0,.35em) -- (.45em,0.6em) -- (.45em,0.1em)  -- cycle (0,.0em) -- (0.5em,.0em);
    }}
% 正规子群
\newcommand{\nsg}{    
    \mathord{
        \tikz[baseline=0pt, trim left=-0.2em, trim right=0.8em] \draw (0,.35em) -- (.45em,0.6em) -- (.45em,0.1em)  -- cycle (0,.35em) -- (0.45em,.35em);
    }}
\newcommand{\nsgeq}{    
    \mathord{
        \tikz[baseline=0pt, trim left=-0.2em, trim right=0.8em] \draw (0,.35em) -- (.45em,0.6em) -- (.45em,0.1em)  -- cycle (0,.35em) -- (0.45em,.35em) (0,.0em) -- (0.5em,.0em);
    }}